% This is samplepaper.tex, a sample chapter demonstrating the
% LLNCS macro package for Springer Computer Science proceedings;
% Version 2.21 of 2022/01/12
%
\documentclass[runningheads]{llncs}
%
\usepackage{natbib}
\usepackage[T1]{fontenc}
\usepackage{amsfonts} 
\usepackage{amsmath}
\usepackage{algorithm}
\usepackage{algorithmic}
\usepackage{xspace}
\usepackage{multirow}
%\usepackage[table]{xcolor}
\usepackage{array}
\usepackage{tabularx}
\usepackage{booktabs} 
\usepackage[table,xcdraw]{xcolor}
\usepackage{subcaption}
\usepackage{amsfonts} 
\usepackage{amsmath}
\usepackage{bm}
\usepackage{bbm}
\usepackage{newfloat}
\usepackage{listings}
\usepackage[draft]{hyperref}
\usepackage{nicematrix}
% T1 fonts will be used to generate the final print and online PDFs,
% so please use T1 fonts in your manuscript whenever possible.
% Other font encondings may result in incorrect characters.
%
\usepackage{graphicx}
% Used for displaying a sample figure. If possible, figure files should
% be included in EPS format.
%
% If you use the hyperref package, please uncomment the following two lines
% to display URLs in blue roman font according to Springer's eBook style:
%\usepackage{color}
%\renewcommand\UrlFont{\color{blue}\rmfamily}
%\urlstyle{rm}
%
\begin{document}
%
\title{LLMs-Empowered Retrieval on Text-Attributed Graphs}
%
%\titlerunning{Abbreviated paper title}
% If the paper title is too long for the running head, you can set
% an abbreviated paper title here
%
\author{Yifei Huang\inst{1}\orcidID{0000-1111-2222-3333} \and
Mingzhe Liu\inst{1}\orcidID{0000-1111-2222-3333}\and
Leilei Sun\inst{1}\orcidID{0000-1111-2222-3333}
\and
Tongyu Zhu\inst{1}\orcidID{0000-1111-2222-3333}
\and
Weifeng Lv\inst{1}\orcidID{0000-1111-2222-3333}
}
%
\authorrunning{Yifei Huang et al.}
% First names are abbreviated in the running head.
% If there are more than two authors, 'et al.' is used.
%
\institute{Beihang University, Beijing 100191, China
\email{\{yifeihuang,mzliu1997,leileisun,zhuyongyu,lwf\}@buaa.edu.cn}}

\maketitle              % typeset the header of the contribution
%
\begin{abstract}
With rich text descriptions and structural information, text-attributed graphs (TAGs) are widely utilized across various applications such as academic networks, social networks, recommendation systems, and knowledge engines. 
In these applications, the task of information retrieval is pivotal for utilizing and interacting with TAGs.
Considering the remarkable and versatile capabilities of large language models (LLMs), we investigate into unleashing LLMs for retrieval on TAGs.
However, it is still an open research question in understanding and leveraging structural patterns with LLMs for effective search for graph nodes and structures.
In this paper, we propose an LLMs-empowered architecture for TAG retrieval named Query Subgraph Schema Guided Retrieval (QSS-Ret).
Our method begins with the generation of a query subgraph schema by LLMs, serving to distill and identify key information from the retrieval query in a graphical form. 
Based on this schema, we retrieve query-relevant subgraphs by locating candidate nodes and adaptable structures under a Depth First Strategy (DFS)-based strategy.
To attend to the fine-grained textual information associated with graph nodes, the extracted subgraphs are summarized and then consumed by the LLM for the re-ranking of candidate graph nodes.
Empirical evaluations demonstrate that our method effectively enhances structural understanding capabilities across TAG retrieval benchmarks from different domains.

\keywords{Text-Attributed Graph  \and Retrieval \and Large Language Model.}
\end{abstract}
%
%
%
\input{texfiles/intro}

\input{texfiles/preliminary}

\input{texfiles/methodology}

\input{texfiles/experiments}

\input{texfiles/related}

\input{texfiles/conclusion}

\section{Acknowledgement}
The authors would like to thank the anonymous reviewers for their constructive comments on this research. This work was supported by the National Natural Science Foundation of China (xxxxxxxx).
%
% ---- Bibliography ----
%
% BibTeX users should specify bibliography style 'splncs04'.
% References will then be sorted and formatted in the correct style.
%
% \bibliographystyle{splncs04}
% \bibliography{mybibliography}
\bibliographystyle{splncs04}
\bibliography{iconip25}
\end{document}
